\begin{agradecimentos}

\subsection*\textit{Por Breno}:

Agradeço, com toda sinceridade, aos meus pais, por sempre colocarem minha educação em primeiro lugar e, com isso, me darem a oportunidade de estudar nas melhores instituições de ensino. Vocês foram excepcionais na arte de criar e educar um filho, mesmo diante de todas as adversidades e dificuldades. Sou grato também à minha tia Lúcia e ao meu tio Gilson, por me acolherem e cuidarem de mim como o próprio filho. À tia Lúcia, um agradecimento especial por ter estudado comigo desde o colégio até o término da faculdade. Agradeço à minha prima Helena, por estar sempre presente nas conversas em momentos de insegurança e dificuldade, e à minha sobrinha Estela, que alegra qualquer ambiente por onde passa. Sou imensamente grato à Cecília, minha namorada, que me motivou constantemente a melhorar e a continuar crescendo como profissional.

Agradeço a todos os professores que nos instruíram da melhor forma ao longo dessa jornada de formação. Deixo um agradecimento especial a dois deles: à Professora Vivian, nossa orientadora neste trabalho, que ao longo de vários meses nos guiou de forma exímia e exemplar; e ao professor Vinícius Gusmão, o melhor professor que já tive, que, mesmo nas matérias mais complexas, fazia de tudo para explicar o conteúdo da forma mais simples possível aos seus alunos, e que, além disso, foi um excelente adversário no tênis de mesa, mesmo com o tempo curto. Por fim, um agradecimento a todos os amigos que fiz na faculdade: Ricardo, minha dupla de TCC e companheiro de graduação; Vinícius, meu baixinho favorito; Pedro, meu companheiro de almoço; e Nathan, Marcos e Alek, com quem passei horas em ótimas conversas. Muito obrigado a todos. 


\subsection*\textit{Por Ricardo}:

Com certeza, este é um momento muito especial. Gostaria de dedicar algumas palavras àqueles que me apoiaram e ofereceram suporte durante esta longa jornada. Em primeiro lugar, gostaria de agradecer à minha família. Meus pais, Ricardo e Maria Beatriz, que me deram a possibilidade de estar onde estou, e que sempre se atentaram ao meu desenvolvimento com todo o carinho do mundo. Minha melhor amiga e irmã, Maria Elisa, que, desde os meus 8 anos, é a minha maior companheira, e meu maior orgulho. Meus avós, Lucrécio, Georgina, Márcio e Marina, que me agraciam com amor incondicional desde que nasci. Minha namorada, Cleo, que me inspira e motiva todos os dias a conquistar tudo que sonho e mereço. E a todos os meus amigos com quem posso contar para tornar a vida mais leve: Pedro, João, Murilo, Lucas, Áxel, Pedro, Batista, Felipe, João e Patrick, sou muito sortudo por ter a amizade de vocês.

Agradeço também à Universidade Federal do Rio de Janeiro, que possibilitou muitos desafios, e com eles, muito aprendizado. Agradeço aos professores que dividiram conhecimento comigo, em especial à professora Vivian, nossa orientadora, que sempre apoiou nossas ideias e nos guiou por todo o processo deste trabalho. E à professora Maria Luiza, minha orientadora acadêmica, que me deu a oportunidade de ajudá-la com a gestão do Minerv@s Digitais, projeto de extensão do qual tenho orgulho imenso de ter participado. E, por fim, agradeço à minha dupla de TCC, Breno. Estivemos lado a lado por boa parte da nossa graduação. Ver este capítulo se encerrando é emocionante, mas não triste, pois tenho a certeza de que, mesmo não fazendo mais parte da rotina um do outro, continuaremos nos incentivando a atingir novos horizontes.

% Os agradecimentos devem ser dirigidos àqueles que contribuíram de maneira relevante à elaboração do trabalho,
% restringindo-se ao mínimo necessário, como instituições (CNPq, CAPES, UFRJ, empresas ou organizações que fizeram parte
% da pesquisa), ou pessoas (profissionais, pesquisadores, orientadores, etc.). 

% Os agradecimentos devem ser colocados de forma hierárquica de importância e para trabalhos financiados com recursos de
% instituições (CAPES, CNPq, FINEP, FAPERJ, etc.) os agradecimentos são obrigatórios a essas instituições.
\end{agradecimentos}